% !TeX root = ../main.tex
\onecolumn
\setcounter{figure}{0}
\setcounter{table}{0}
\renewcommand{\thefigure}{A\arabic{figure}}
\renewcommand{\thetable}{A\arabic{table}}
\section{Corpus and Platform Dictionaries}
\label{app:corpus}
Figure~\ref{fig:corpus} reports indexed contributions by edition. These are publication counts, not submission or acceptance counts. The main analysis includes all retained main-conference formats.
\begin{figure}[ht]
\centering
\includegraphics[width=.75\textwidth]{figures/01_corpus.png}
\caption{Annual contributions in the retrieved corpus.}
\label{fig:corpus}
\end{figure}
\subsection{Platform families}
\label{app:platforms}
The twenty named platform families are Twitter/X, Reddit, Facebook, Wikipedia, YouTube, Instagram, TikTok/Douyin, Weibo, WeChat, WhatsApp, Telegram, Gab, 4chan, Parler, Mastodon, Bluesky, Flickr, Digg, MySpace, and StackOverflow. A separate weblog indicator matches blogs, weblogs, blogging, bloggers, and blogosphere. Twitter includes tweets and retweets; Reddit includes subreddit and redditor variants; Wikipedia includes Wikipedian variants. StackOverflow allows a space between the two words. Seven X-only cases, identified through AI-assisted inspection, supplement the Twitter vocabulary. All other named families use the explicit name, with case-insensitive word boundaries. The regular expressions used in the analysis will be released upon publication and provide the exact specification. Rates count mentions, can overlap, and do not establish data use.

\section{Topic Components and Complete Distributions}
\label{app:topics}
Topic labels describe one exploratory 15-component NMF representation. Table~\ref{tab:topicterms} lists each component's twenty highest-weight terms in order. Figure~\ref{fig:topicsproblems} orders rows by the difference between the last and first fixed periods, from the largest increase to the largest decrease, separately for topics and framings. Component weights and cue prevalences have different denominators.
\begingroup
\small
\begin{longtable}{@{}p{.05\textwidth}p{.22\textwidth}p{.65\textwidth}@{}}
\caption{NMF components and top terms, ordered by component identifier.}\label{tab:topicterms}\\
\toprule
ID & Interpretive label & Twenty highest-weight terms \\
\midrule
\endfirsthead
\toprule
ID & Interpretive label & Twenty highest-weight terms \\
\midrule
\endhead
\bottomrule
\endfoot
T00 & Networks, diffusion, recommendation & diffusion, influence, link, graph, structure, recommendation, nodes, problem, prediction, links, cascades, node, systems, graphs, real, items, spread, popularity, link prediction, algorithms \\ \addlinespace
T01 & Conventional classification & learning, classification, features, detection, machine, machine learning, performance, task, accuracy, supervised, art, state art, training, state, framework, labeled, reviews, f1, problem, existing \\ \addlinespace
T02 & News and journalism & news, articles, news articles, article, outlets, sources, news outlets, fake, news sources, coverage, fake news, stories, news article, comments, misinformation, bias, source, readers, sharing, local news \\ \addlinespace
T03 & Language modeling and generative AI & llms, language, llm, ai, human, language llms, generated, gpt, tasks, fine, performance, text, reasoning, english, languages, framework, bias, biases, natural language, detection \\ \addlinespace
T04 & Pandemics and vaccination & covid, covid 19, 19, pandemic, 19 pandemic, vaccine, anti, misinformation, vaccination, vaccines, related, impact, health, stance, concerns, narratives, public, understand, population, conspiracy \\ \addlinespace
T05 & Politics and polarization & political, election, accounts, elections, polarization, discourse, party, politicians, partisan, campaigns, presidential, public, parties, leaning, engagement, right, political discourse, campaign, topics, candidates \\ \addlinespace
T06 & Sentiment and opinion mining & sentiment, topic, topics, text, sentiments, opinion, positive, negative, opinions, corpus, words, polarity, word, set, mood, semantic, emotions, related, mining, terms \\ \addlinespace
T07 & Communities and participation & communities, community, members, group, groups, participation, activity, subreddits, behavior, moderation, work, engagement, discussion, interactions, active, patterns, success, subreddit, dynamics, effect \\ \addlinespace
T08 & Events and collective attention & events, event, real, world, real world, time, real time, world events, messages, stream, response, summarization, event detection, crisis, temporal, streams, disaster, detection, short, relevant \\ \addlinespace
T09 & Identity and self-presentation & gender, profiles, people, profile, personality, self, age, friends, differences, personal, female, women, attributes, men, traits, male, likely, demographic, interactions, groups \\ \addlinespace
T10 & Search, Q\&A, and tagging & question, search, questions, answering, answer, answers, quality, question answering, tagging, tags, systems, queries, engine, search engine, community question, engines, search engines, knowledge, expertise, tag \\ \addlinespace
T11 & Health and social support & health, mental, mental health, support, posts, self, individuals, depression, emotional, public, post, risk, effects, findings, stress, public health, related, treatment, examine, discuss \\ \addlinespace
T12 & Hateful and abusive speech & hate, speech, hate speech, speech detection, detection, hateful, moderation, targets, target, attacks, comments, replies, toxic, offensive, understanding, harassment, automatic, language, conversations, toxicity \\ \addlinespace
T13 & Visual and multimodal media & videos, video, images, memes, visual, meme, image, multimodal, popularity, misinformation, channels, sharing, creators, popular, whatsapp, shared, channel, metadata, short, audio \\ \addlinespace
T14 & Location and mobility & location, mobility, urban, check, patterns, foursquare, city, locations, temporal, services, cities, spatial, places, geographic, human mobility, ins, check ins, behavior, time, mobile \\ \addlinespace
\end{longtable}
\endgroup

\begin{figure}[p]
\centering
\includegraphics[width=\textwidth]{figures/04_topics_problems.png}
\caption{Complete annual topic shares and framing-cue prevalences, ordered by endpoint-period change. Topic shares sum to 100\%; framing categories overlap. No historical phase boundaries are imposed.}
\label{fig:topicsproblems}
\end{figure}
\section{Framing Codebook}
\label{app:codebook}
The definitions distinguish an intended research problem from the lexical rule used to approximate it. Example cues are illustrative; the regular expressions used in the analysis specify matching exactly. A cue can be incidental, and an absent cue does not establish an absent concern. Categories are not mutually exclusive.
\begingroup
\small
\begin{longtable}{@{}p{.20\textwidth}p{.46\textwidth}p{.26\textwidth}@{}}
\caption{Definitions and example lexical cues for the eight overlapping framing cues.}\label{tab:codebook}\\
\toprule
Category & Definition and scope limitation & Example cues \\
\midrule
\endfirsthead
\toprule
Category & Definition and scope limitation & Example cues \\
\midrule
\endhead
\bottomrule
\endfoot
P1 Structure and diffusion & Explicit framing of relational structure, diffusion, contagion, social influence, collective dynamics, or community formation. Does not identify every explanatory study of collective behavior. & diffusion; cascades; social ties; homophily \\ \addlinespace
P2 Prediction and classification & Explicit prediction, forecasting, classification, detection, or recommendation terminology. May match a prediction or detection task that a paper critiques rather than performs. & prediction; forecasting; classification; detection \\ \addlinespace
P3 Measurement and validity & Explicit validation, sampling, representation, reliability, ground truth, reproducibility, or measurement-error terminology. Routine uses of the verb measure alone do not qualify. Validation and reliability can refer to routine evaluation or background discussion, not a central measurement contribution. & validity; sampling; ground truth; reliability \\ \addlinespace
P4 Explanation and intervention & Explicit causation, intervention, recognizable experimental-design terminology, causal/social mechanisms, or theory testing. Generic mechanism and theoretical implication mentions are excluded. Theory, mechanism, or causal language does not establish a valid causal design. & causal; intervention; randomized; propensity score \\ \addlinespace
P5 Harm and information integrity & Explicit information-integrity, abuse, discrimination, privacy-risk, spam, or other harm framing. Includes descriptive and technical work about harm, without inferring the authors' normative commitments. & misinformation; spam; harassment; privacy; harm \\ \addlinespace
P6 Governance and participation & Explicit moderation, governance, censorship, enforcement, platform rules, or community norms. Does not capture every study of participation; moderation can also be mentioned as background. & moderation; governance; censorship; community rules \\ \addlinespace
P7 Support and well-being & Explicit human-health, care, distress, social-support, or well-being framing. The bare word health is excluded to reduce organizational-health metaphors. A health mention does not establish a clinical outcome or a supportive intervention. & mental health; caregiving; depression; social support \\ \addlinespace
P8 Resources and research access & Explicit nearby provision of data/tools or data-access and research-infrastructure terminology. Generic API and publicly available mentions alone do not qualify. A paper can provide a tool without belonging to a dataset or demonstration track. & researcher access; open-source; releases a dataset \\ \addlinespace
\end{longtable}
\endgroup

\section{Sensitivity Results}
\label{app:sensitivity}
Table~\ref{tab:sensitivity} reports the harm and community-governance checks using the four main comparison periods. Density is the mean of paper-level cue matches divided by title/abstract word count and multiplied by 100; it is not a paper prevalence. The historical union adds trolling, flaming, and vandalism terms to all editions. Spam is already in the baseline harm lexicon. The strict governance cue excludes papers whose only matches are \emph{moderate}, \emph{moderated}, \emph{moderates}, or \emph{moderating}.
\begingroup
\small
\begin{longtable}{@{}p{.36\textwidth}rrrr@{}}
\caption{Sensitivity estimates. All rows except cue density are percentages.}\label{tab:sensitivity}\\
\toprule
Measure & 2007--11 & 2012--16 & 2017--21 & 2022--26 \\
\midrule
\endfirsthead
\toprule
Measure & 2007--11 & 2012--16 & 2017--21 & 2022--26 \\
\midrule
\endhead
\bottomrule
\endfoot
Harm: full title/abstract & 6.2 & 12.3 & 27.0 & 38.1 \\ \addlinespace
Harm: title + first 100 abstract words & 5.7 & 9.5 & 22.7 & 31.9 \\ \addlinespace
Harm: first 100 abstract words only & 5.7 & 8.7 & 21.9 & 31.0 \\ \addlinespace
Harm: historical union & 6.2 & 12.3 & 27.9 & 38.3 \\ \addlinespace
Harm: mean matches / 100 words & 0.10 & 0.23 & 0.54 & 0.78 \\ \addlinespace
Governance within community topic & 4.0 & 1.2 & 14.4 & 34.3 \\ \addlinespace
Governance after feature masking & 3.6 & 1.1 & 13.4 & 32.5 \\ \addlinespace
Governance within community: title + first 100 words & 2.5 & 0.8 & 8.8 & 25.8 \\ \addlinespace
Governance within community: first 100 abstract words & 2.5 & 0.8 & 7.4 & 25.6 \\ \addlinespace
Governance within community: strict cue & 2.5 & 0.5 & 12.5 & 33.3 \\ \addlinespace
Governance cues overall: strict cue & 0.8 & 0.6 & 5.5 & 13.1 \\ \addlinespace
\end{longtable}
\endgroup

\subsection{Direct term overlap}
Table~\ref{tab:overlap} gives the number of each topic's top twenty terms matched by each framing expression. For the community component, only \emph{moderation} matches the governance expression. The vocabulary-removal analysis masks every governance-matching unigram or bigram in the fitted vocabulary, not just the top twenty terms. The complete term-level overlap record will be released with the analysis materials upon publication.
\begingroup
\small
\begin{longtable}{@{}lrrrrrrrr@{}}
\caption{Counts of top-twenty topic terms matched by each framing lexicon. These are literal term matches, not measures of semantic independence.}\label{tab:overlap}\\
\toprule
Topic & P1 & P2 & P3 & P4 & P5 & P6 & P7 & P8 \\
\midrule
\endfirsthead
\toprule
Topic & P1 & P2 & P3 & P4 & P5 & P6 & P7 & P8 \\
\midrule
\endhead
\bottomrule
\endfoot
T00 & 2 & 3 & 0 & 0 & 0 & 0 & 0 & 0 \\ \addlinespace
T01 & 0 & 2 & 0 & 0 & 0 & 0 & 0 & 0 \\ \addlinespace
T02 & 0 & 0 & 0 & 0 & 2 & 0 & 0 & 0 \\ \addlinespace
T03 & 0 & 1 & 0 & 0 & 0 & 0 & 0 & 0 \\ \addlinespace
T04 & 0 & 0 & 0 & 0 & 2 & 0 & 7 & 0 \\ \addlinespace
T05 & 0 & 0 & 0 & 0 & 0 & 0 & 0 & 0 \\ \addlinespace
T06 & 0 & 0 & 0 & 0 & 0 & 0 & 0 & 0 \\ \addlinespace
T07 & 0 & 0 & 0 & 0 & 0 & 1 & 0 & 0 \\ \addlinespace
T08 & 0 & 2 & 0 & 0 & 0 & 0 & 0 & 0 \\ \addlinespace
T09 & 0 & 0 & 0 & 0 & 0 & 0 & 0 & 0 \\ \addlinespace
T10 & 0 & 0 & 0 & 0 & 0 & 0 & 0 & 0 \\ \addlinespace
T11 & 0 & 0 & 0 & 0 & 0 & 0 & 3 & 0 \\ \addlinespace
T12 & 0 & 2 & 0 & 0 & 5 & 1 & 0 & 0 \\ \addlinespace
T13 & 0 & 0 & 0 & 0 & 1 & 0 & 0 & 0 \\ \addlinespace
T14 & 0 & 0 & 0 & 0 & 0 & 0 & 0 & 0 \\ \addlinespace
\end{longtable}
\endgroup

\section{Exploratory Temporal Segmentation}
\label{app:boundaries}
For segmentation only, we normalize annual overlapping framing-cue rates together with an unmatched category. We concatenate square-root topic and framing shares, weighting the axes equally, and minimize within-segment squared deviation. Each edition has equal weight in this diagnostic. This differs from pooling papers within the fixed periods used in the main analysis.

The baseline five-segment solution with a three-edition minimum has starts in 2007, 2011, 2017, 2021, and 2024. Its full partition recurs in 29/200 resamples (14.5\%). Exact boundary frequencies are given below. The baseline minimum makes 2024 the latest possible final start, so its apparent recurrence cannot establish a recent historical break. A second analysis uses the same 200 newly sampled paper compositions to compare two- and three-edition minima. It admits a 2025 start, which appears in 47/200 resamples (23.5\%). These are conditional exact-year frequencies, not confidence levels for historical events.
\begingroup
\small
\begin{longtable}{@{}lrrr@{}}
\caption{Exact-year boundary recurrence (percent) under five-segment joint profiles. A dash marks a boundary disallowed by the minimum-length rule. The two new columns use identical resampled paper compositions.}\label{tab:boundaries}\\
\toprule
Boundary start & Baseline 200; min. 3 & New 200; min. 3 & Same new 200; min. 2 \\
\midrule
\endfirsthead
\toprule
Boundary start & Baseline 200; min. 3 & New 200; min. 3 & Same new 200; min. 2 \\
\midrule
\endhead
\bottomrule
\endfoot
2009 & -- & -- & 21.5 \\ \addlinespace
2010 & 39.5 & 39.0 & 33.0 \\ \addlinespace
2011 & 56.5 & 59.0 & 54.5 \\ \addlinespace
2012 & 3.5 & 1.5 & 28.0 \\ \addlinespace
2013 & 19.5 & 20.5 & 5.0 \\ \addlinespace
2014 & 16.5 & 13.0 & 8.0 \\ \addlinespace
2015 & 6.5 & 9.0 & 5.5 \\ \addlinespace
2016 & 1.0 & 2.5 & 0.5 \\ \addlinespace
2017 & 54.0 & 53.0 & 55.5 \\ \addlinespace
2018 & 34.0 & 34.0 & 31.5 \\ \addlinespace
2019 & 11.5 & 12.5 & 23.0 \\ \addlinespace
2020 & 21.5 & 19.5 & 14.5 \\ \addlinespace
2021 & 65.0 & 65.5 & 70.5 \\ \addlinespace
2022 & 7.0 & 5.5 & 4.0 \\ \addlinespace
2023 & 0.0 & 0.0 & 0.0 \\ \addlinespace
2024 & 64.0 & 65.5 & 21.5 \\ \addlinespace
2025 & -- & -- & 23.5 \\ \addlinespace
\end{longtable}
\endgroup

We also retain results for three, four, and five segments and topic-only, framing-only, and joint representations. The segment count is not selected by an external validation criterion. Changes in component count and initialization can move boundaries. We therefore do not name a definitive sequence of five phases or use it to calculate the headline period comparisons. Broad emphases can overlap even when a segmentation algorithm must assign each year to one interval.

\section{Independent Validation Protocol}
\label{app:validation}
The framing-validation protocol samples 20 papers per edition (400 total), without restricting to cue-positive cases. Two coders independently judge whether each abstract expresses a substantive objective matching each category, distinguishing central or evaluated concerns from incidental background. Multiple categories are allowed; uncertain cases receive a separate flag. The protocol materials withhold automatic predictions. Source links can reveal bibliographic identity; an optional local reader also withholds author and edition information.

Report per-category Cohen's $\kappa$, raw agreement, and positive-label counts overall and by fixed period. After preserving the independent ratings, adjudicate disagreements and uncertainty. Compare the frozen cues with those adjudicated judgments, reporting precision, recall, and the TP/FP/FN/TN counts for each category and period. Since each edition contributes twenty coded papers regardless of its size, pooled period estimates should use inverse-inclusion weights $N_t/20$. Report denominators and unresolved judgments; undefined recall for a category with no positive gold labels must not be replaced by zero. Development-sample overlaps are flagged for a separate sensitivity analysis. Any lexicon revision based on these judgments requires a new held-out evaluation.

For topic interpretability, the protocol contains five word-intrusion items per component (75 total). Each presents five high-weight terms and an intruder drawn from another component, with an IDF-matching rule to reduce rarity cues. Eighty topic-intrusion items (four papers per edition) present the three highest-weight topic word lists and one low-weight list for the document. Coders select the intruder without seeing the model weights or key. Report correct-choice rates, counts, uncertainty, and comments by component or period; chance rates are $1/6$ and $1/4$, respectively. These tasks assess interpretability; they do not independently establish the substantive validity of a label.
